\subsection{Design}

With the deterministic grounding layer of \S\ref{sec:grounding-layer} in
place, we ran a 120-trial real-hardware study comparing a static-scene LLM
planner against a camera-grounded VLM planner, methodologically anchored
to the task/environment complexity taxonomy and the TS/TSR/AETS metric
family of Favali et al.~\cite{favaliRAS}, adapted to a physical rig. Six
task conditions crossed three difficulty axes, task length, reference
specificity, and affordance/perception load, each at an easy and a hard
extreme, over a real scene of four color-distinguished cubes and a tray:
\code{pp\_easy}/\code{pp\_hard} (pick-and-place, one vs.\ four objects,
named directly), \code{sort\_easy}/\code{sort\_hard} (stacking and
sorting), and \code{arith\_easy}/\code{arith\_hard} (the object to act on
must be \emph{computed} from a stated numeric value mapping rather than
named directly, the one condition designed to isolate reasoning load
from perception load). Each condition was run 10 times per model, for
120 trials total, with a single reasoning method (chain-of-thought with
self-consistency) held fixed across all trials, since it is reported as
the most safety-robust strategy in~\cite{favaliRAS} and avoids stacking
reasoning-method variance on top of the LLM-vs-VLM comparison actually
under test. Both the LLM and VLM planner use the same underlying model
(\code{kimi-k2.6}) for text-only and image-grounded calls respectively.
This is deliberate, so that grounding modality is the only variable
isolated between arms, not a side effect of comparing two different
models with different baseline capabilities.\footnote{An earlier run used
different models per modality; the VLM-specific model was subsequently
withdrawn from its provider's catalog, forcing a full re-run on a single
unified model instead.}

\subsection{Metrics, adapted for physical ground truth}

Favali et al.'s TS (task safety), TSR (task success rate), and AETS
(action efficiency on task success) metrics are fully automatic in a
symbolic simulator, where whether an action ``damages the environment''
or a sub-task ``completes'' is always exactly derivable from simulator
state. Our real UR5cb has no collision or force-fault sensor anywhere in
its control path, and no vision-based verification of the final scene
state. We therefore adapt the same three metrics to a
\textbf{semi-automatic} protocol: everything the system can know
automatically (the command, configuration, generated plan, and steps
actually executed) is logged by the debug harness of
\S\ref{sec:grounding-layer} for every trial; the operator supplies exactly
two numbers immediately after each trial, was it safe, and how many
sub-tasks completed correctly, from which TS, TSR, and AETS are
computed identically to~\cite{favaliRAS}'s definitions. This adaptation is
itself a consequence of the paper's central argument: full automation of
these metrics is only possible because a symbolic simulator's world state
is always exactly known, which is precisely the property real hardware
does not have.

At $n=10$ per cell, point estimates alone invite exactly the kind of
pushback a reviewer should raise, so every reported TS\% and TSR\% below
carries a 95\% Wilson score interval (TS is an exact per-trial
proportion; TSR pools successes over $\text{sub\_tasks\_required} \times
n$). One comparison was declared before any data was collected, not
chosen afterward to fit the numbers: a two-sided Fisher's exact test,
LLM vs.\ VLM, on \code{arith\_hard} full-trial success, at $\alpha=0.05$.
It is binarized per trial rather than pooled across a trial's four
sub-tasks, since those are not independent draws; one wrong decision
tends to misplace every object in the trial together, not each
independently. Everything else is reported descriptively only.

\subsection{Results}

\begin{table}[t]
\centering
\caption{Overall results, 60 trials per model (120 total), with the
deterministic grounding layer of \S\ref{sec:grounding-layer} in place.}
\label{tab:overall}
\begin{tabular}{@{}lcccc@{}}
\toprule
\textbf{Model} & \textbf{n} & \textbf{TS} & \textbf{TSR} & \textbf{AETS} \\
\midrule
LLM (\code{kimi-k2.6}) & 60 & 98.3\% & 98.3\% & 0.1269 \\
VLM (\code{kimi-k2.6}) & 60 & 85.0\% & 89.2\% & 0.1149 \\
\bottomrule
\end{tabular}
\end{table}

\begin{figure}[t]
\centering
\includegraphics[width=0.95\linewidth]{overall_summary.png}
\caption{Overall TS/TSR/AETS summary, LLM vs.\ VLM, with the deterministic
grounding layer in place.}
\label{fig:overall}
\end{figure}

Table~\ref{tab:overall} reports the headline result. Both models operate
under the same deterministic grounding layer and the same underlying
model, so the residual gap is neither a grasp-geometry gap
(\S\ref{sec:grounding-layer} applies uniformly to both pipelines) nor a
model-capability gap. It is specifically attributable to grounding
modality. Breaking results down by task condition (Figure~\ref{fig:by-task})
shows the gap concentrated on the two hardest multi-object conditions,
and predominantly in \emph{safety}, not success rate: TS is 100\% for
both models on every easy condition and on \code{pp\_hard}/\code{sort\_hard}
(LLM), but drops to 50\% (VLM, \code{sort\_hard}) and 60\%/90\% (VLM/LLM,
\code{arith\_hard}) on the two conditions with the most simultaneous
in-zone objects. Neither model shows the sharp \emph{success-rate}
collapse a July pilot with different, modality-specific models had found
on the arithmetic conditions specifically (30\%/25\% VLM TSR against
100\%/97\% LLM). With a single unified model, \code{arith\_easy} is
100\%/100\% and \code{arith\_hard} is 90\%/85\% (TSR, LLM/VLM): a real but
modest gap, not a collapse. \S\ref{sec:computation-not-perception}
returns to why. The one pre-registered comparison, Fisher's exact test,
LLM vs.\ VLM, \code{arith\_hard} full-trial success, shows no detectable
difference at this sample size (LLM 6/10, VLM 4/10, $p=0.66$). That is
consistent with a real but modest gap that $n=10$ per cell cannot
resolve to significance, not with an absence of one.

\begin{figure}[t]
\centering
\includegraphics[width=0.95\linewidth]{ts_tsr_by_task.png}
\caption{Task safety and task success rate by task condition, with 95\%
Wilson CI.}
\label{fig:by-task}
\end{figure}

\subsection{The failure mode is computation, not perception}
\label{sec:computation-not-perception}

Reading all 22 non-empty operator failure notes from the 120 trials,
rather than trusting the aggregate numbers alone, was the key analysis
step, as it was in an earlier pilot on this same task matrix. Every one
of the 22 notes describes either a collision between objects (10) or an
execution/grasp failure, a slipped grasp, a release too high, a dropped
object (12), during a multi-object task (\code{pp\_hard},
\code{sort\_hard}, \code{arith\_hard}). \textbf{Zero} notes describe a
wrong object chosen for reasoning or value-mapping reasons, and zero
describe an incomplete enumeration (Figure~\ref{fig:failure-modes}). This
happens in \textbf{both} arms: the LLM contributes 4 of the 22 notes (all
on \code{arith\_hard}), the VLM the remaining 18, spread across all three
multi-object conditions roughly in proportion to their trial counts. The
failure recurs regardless of which planner grounded the scene, which is
direct confirmation that Case 6 (\S\ref{sec:case-collision}) is this
taxonomy's most consequential open item, not a peculiarity of one
pipeline.

\begin{figure}[t]
\centering
\includegraphics[width=0.85\linewidth]{failure_modes.png}
\caption{Failure-mode breakdown, derived from the 22 non-empty operator
notes out of 120 trials.}
\label{fig:failure-modes}
\end{figure}

Unlike the closed bugs of \S\ref{sec:taxonomy}, this result is not a
second-order instance of the paper's central claim; it is a direct one.
Case 6 is itself one of the six geometry-grounding bugs of
Table~\ref{tab:taxonomy}; this evaluation is what identifies it as the
dominant one once the other five are closed.

\emph{A note on an earlier finding.} A pilot run on this same task matrix,
using different, modality-specific models (not the single \code{kimi-k2.6}
used here), found a different dominant category: 13 of 31 notes read as a
wrong-object choice concentrated specifically in the arithmetic condition,
consistent with an off-by-one error in value-mapping reasoning under a
stated numeric mapping. That finding does not replicate here: zero such
notes this cycle, on every condition. We read the simplest explanation as
the correct one: it was a property of the earlier VLM model specifically,
not of VLM grounding in general, and does not survive a model change. We
report it for the record and as a caution against generalizing a
single-model finding to ``VLM grounding'' as a category, not as a result
that still stands.
